% !TEX root = ../Dissertation.tex

\chapter{Einleitung}
\label{Einleitung}
\dictum[{\citep[43]{Labov.1984}}]{At the heart of social and emotional expression is the linguistic feature of intensity.}
\smallskip

\dictum[{\citep[110]{Ayren.1986}}]{Sprache lebt, und was lebt, wandelt sich.}


\section{Phänomen und Untersuchungsgegenstand}\label{abs:problem}
Sprachliche Elemente neigen dazu, sich abzunutzen und ihre Leuchtkraft zu verlieren, wenn sie zu häufig gebraucht werden. Hierzu gehören bspw. Metaphern, die -- wie selbst in der Belletristik angemerkt wird {\citep[vgl.][453f.]{Mercier.2020}} -- zu einem pragmatisch bedingten Sprachwandel neigen. Die vorliegende Studie untersucht vor diesem Hintergrund Intensitätspartikeln, die im Deutschen nicht bloß in vielerlei Gestalt auftreten, sondern auch den Ruf haben, besonders anfällig für sprachlichen Wandel -- gar \glqq{}Modewörter\grqq{} \citep[267]{Becher.1907} -- zu sein. Letzteres gilt als ausschlaggebend dafür, dass es eine Unmenge an Ausdrücken gibt, auf die wir als Sprecher:innen\footnote{Im Buch wird im Falle von realen Personen auf genderneutrale Formen (bspw. \textit{Muttersprachler:innen} oder \textit{Versuchspersonen}) zurückgegriffen, bei linguistischen Standardbezeichnungen (\textit{der Sprecher} oder \textit{der Adressat}) das jeweilige grammatische Geschlecht verwendet.} nach Belieben zurückgreifen können
(vgl. \citealt[74]{Baumgarten.1908}; \citealt[379]{Mendez-Naya.2003}). So können wir bspw. sagen, dass der Strand \textit{ziemlich schön} ist, aber er kann auch \textit{sehr schön} sein. Doch nicht nur das: Wenn uns diese Formulierung nicht ausdrucksstark genug ist, können wir ihn auch als \textit{mega schön}, \textit{sau schön} oder \textit{super schön} bezeichnen. Der Literatur nach zeigen wir mithilfe dieser Ausdrucksweisen an, dass die durch das Bezugswort beschriebene Eigenschaft (d. h. die Schönheit) nicht nur in einem normalen Ausmaß vorliegt, sondern besonders stark ausgeprägt ist (vgl. {\citealt[22]{Paradis.1997}}; {\citealt[133]{Gutzmann.2019a}}). In diesem Fall würde sie auf einer zugehörigen Skala einen überdurchschnittlich hohen Wert erreichen, der eine mögliche durch den Bezugsausdruck \textit{Strand} determinierte Norm übersteigt. Daneben bringen wir mit der Äußerung den Annahmen von {\citet[315]{Lang.1983}} und {\citet[57]{Ortner.2014}} zufolge ebenso unsere persönliche Einstellung zum besagten Sachverhalt zum Ausdruck, die sich wiederum in verschiedenen affektiven Qualitäten wie Begeisterung, Überraschung oder Überzeugung niederschlagen kann. Sprachliche Formulierungen, die mit einem solch hohen Werte- und Einstellungsgrad einhergehen, nennen sich Intensivierungen. Darunter versteht man im Wesentlichen die \glqq{}Verstärkung des emotionalen Gehalts eines Ausdrucks und/oder der Intensität einer bezeichneten Eigenschaft durch formale Mittel.\grqq{} ({\citealt[299]{Glueck.2000}}). Die dafür zur Verfügung stehenden Lexeme werden gemeinhin als Intensitäts- oder Intensivpartikeln -- kurz Intensivierer (englisch: \textit{intensifiers}) -- bezeichnet.

%\footnote{In der Literatur werden die Begriffe zumeist synonym verwendet. Um mögliche Unklarheiten zu vermeiden, werden in dieser Arbeit ausschließlich 'Intensitätspartikeln' und 'Intensivierer' verwendet; diese jedoch vollkommen bedeutungsgleich.}


%\begin{singlespace*}
%\begin{quotation} "Since adverbs have the effect of intensifying the adjectives they modify, they are used to give specifications of degree and moreover they show involvement on the part of the speaker. In this respect, they add to the emotive and subjective dimension of discourse" \citep[555]{Athanasiadou.2007}. \end{quotation}
%\end{singlespace*}
Gegenstand dieser Untersuchung ist das in der theoretischen Literatur unlängst stark diskutierte Phänomen der lexikalischen Adjektiv-Intensivierung, dessen Erforschung in der Sprachwissenschaft in den letzten dreißig Jahren u. a. durch die Arbeit von \citet[]{Os.1989} erheblich an Dynamik gewonnen hat. Dabei wird ein graduierbares Adjektiv durch eine Intensitätspartikel modifiziert. Infolge der Modifikation erfährt das Adjektiv eine Bedeutungsverstärkung, die sich in einer gesteigerten Intensität und einer höheren Position auf der zugehörigen (mentalen) Skala manifestiert. Da vor allem Adjektive von sprachlicher Intensivierung betroffen sind, haben sich empirische soziolinguistische Untersuchungen in diesem Bereich bisher hauptsächlich auf die Intensivierung von Adjektivköpfen konzentriert (vgl. bspw. \citealt[]{Stratton.2020b}).


\section{Deskriptive versus expressive Intensitätspartikeln}
\label{abs:deexin}

Bei einem Blick in die einschlägige Literatur zu Intensitätspartikeln fällt auf, dass diese grundsätzlich in zwei Klassen unterteilt werden können: \glqq{}deskriptive Intensivierer\grqq{} \citep[]{Gutzmann.2012} und \glqq{}expressive Intensivierer\grqq{} (ebd.). Die Klasse der deskriptiven Intensivierer -- zu denen Ausdrücke wie \textit{äußerst}, \textit{überaus} oder \textit{sehr} gehören -- gilt als relativ stabil, da es sich um ein fest etabliertes lexikalisches Inventar handelt, dessen Umfang sich kaum wahrnehmbar verändert. In der theoretischen Literatur wird diese Klasse meist durch die deutsche Standard-Intensitätspartikel \textit{sehr} beschrieben:

\ea
\ea Das Hotel ist sehr schön.
\ex Das Essen ist sehr lecker.
\z
\z
Deskriptive Intensivierer gelten als delexikalisiert und damit konnotationslos, weshalb sie eine neutral skalierende Wirkung haben. Ihr semantischer Beitrag liegt ausschließlich darin, den Ausprägungsgrad der durch das Adjektiv ausgesagten Eigenschaft auf einer Skala zu erhöhen.\footnote{Unter \textit{neutral} verstehe ich Ausdrücke, die grundsätzlich keinerlei emotional-evaluative Färbung aufweisen und folglich weder positiv noch negativ konnotiert sind.} Über die Intensitätssteigerung hinaus haben die Ausdrücke der deskriptiven Intensiviererklasse keinen Einfluss auf die Bedeutung des intensivierten Adjektivs.

Die Klasse der expressiven Intensivierer gilt dagegen als deutlich heterogener und weniger klar abgegrenzt als die der deskriptiven Intensivierer; zudem umfasst sie verschiedene Sprachvarietäten, etwa Jugend- oder Vulgärsprache. Beispiele hierfür sind Ausdrücke wie \textit{mega}, \textit{scheiße} oder \textit{verdammt}:


\ea
\ea Der Flug ist mega kurz.
\ex Der Bus ist scheiße voll.
\ex Das Wetter ist verdammt schlecht.
\z
\z
Expressive Intensivierer treten vor allem in der Umgangssprache auf, häufig in morphologisch ähnlicher Gestalt wie Adjektive, und haben eine vergleichbare Funktion wie der klassische deskriptive Intensivierer \textit{sehr}: So wird in beiden Fällen verglichen mit dem Positiv ein höheres Ausmaß einer Eigenschaft ausgedrückt. Der durch expressive Intensivierer angezeigte Grad wird jedoch -- so die in der Literatur vorherrschende Annahme -- intuitiv als stärker empfunden. Darüber hinaus offenbart die Partikel durch eine emotional-expressive Zusatzkomponente zugleich die persönliche Einstellung des Sprechers zum thematisierten Sachverhalt, die nicht Teil des propositionalen Gehalts ist {\citep[vgl.][5]{Gutzmann.2013a}}. Anders als deskriptive Intensivierer verstärken expressive Intensivierer nicht nur das durch das Adjektiv beschriebene Merkmal, sondern spiegeln auch die subjektive Haltung des Sprechers wider. Dadurch werden sie vom Hörer als affektiver wahrgenommen. Aufgrund dieser expressiven Komponente ist ihr Gebrauch stärker restringiert; sie finden sich typischerweise in informellen und mündlichkeitsnahen Sprachregistern -- besonders häufig in der Jugendsprache. Deskriptive Intensivierer hingegen können in nahezu allen Kontexten problemlos verwendet werden. Zusammengefasst versteht man unter expressiven Intensivierern also Ausdrücke, die als einstellungsausdrückend wahrgenommen werden.\footnote{Zu beachten ist, dass deren tatsächlicher Expressivitätsgrad erst durch die in diesem Buch dargelegten experimentellen Untersuchungen ermittelt wird.} Mit ihnen signalisiert der Sprecher zum einen, dass eine im Sachverhalt bezeichnete Eigenschaft über ein als dafür typisch geltendes Durchschnittsmaß hinausgeht, und er dieser Normabweichung zum anderen eine besondere Bedeutung beimisst. Hierbei trägt er die semantische Rolle des Experiencers. Die Gemeinsamkeiten und Unterschiede zwischen deskriptiven und expressiven Intensivierern werden in Abschnitt {\ref{Intensivierertypen}} im Detail untersucht.

%Die Partikeln selbst tragen keine eigenständige Bedeutung, sondern orientieren sich hinsichtlich der Polarität an dem nachfolgenden Adjektiv, das sie modifizieren. Oder mit anderen Worten: "[I]t looks like the lexical content of the adjacent phrase is a good indicator of emotional polarity" \citep[12]{Constant.2009}.

%Gleiches findet sich in \citet[312]{Pfeifer.1995}: 
%
%\begin{singlespace*}
%\begin{quotation} \textbf{expressiv} Adj. 'ausdrucksvoll', im 18. Jh. entlehnt aus gleichbed. frz. \textit{expressif}, einer Ableitung von frz. \textit{expression} 'Ausdruck' \textelp{} Substantivbildung zum Verb \textit{exprimere} 'heraus-, ausdrücken, anschaulich machen'. \end{quotation}
%\end{singlespace*}
%Doch was bedeutet es eigentlich, wenn ein sprachlicher Ausdruck {\textit{expressiv}} ist? Definitionsgemäß wird mit dem aus dem Lateinischen stammenden Adjektiv 'etwas zum Ausdruck' gebracht, was laut dem Duden-Universalwörterbuch im Sinne von 'ausdrucksvoll, ausdrucksstark, ausdrucksbetont, mit Ausdruck' zu verstehen ist {\citep[vgl.][576]{Duden.2019}}. Nach {\citet[193]{Glueck.2000}} handelt es sich dabei um "Textelemente, die aufgrund ihrer stilist. Gestaltung in besonderem Maße Emotionen und Einstellungen des Sprechers ausdrücken". Als subjektive Einstellungsbekundungen liefern Expressiva also Informationen über die inneren Zustände, Einstellungen und Bewertungen des Sprechers, die dem Adressaten andernfalls verborgen blieben. Hierin unterscheiden sie sich von rein neutralen Sachverhaltsbeschreibungen: Expressiva beschreiben innere Zustände nicht, sondern sie bringen sie schlicht freiweg zum Ausdruck {\citep[vgl.][485]{Gutzmann.2018}}. Damit geht offenbar ein höherer Intensitätsgrad einher, der sich nach Auffassung vieler Autor:innen auf eine zugehörige mentale Skala projizieren lässt (vgl. u.\,a. {\citealt[]{Biedermann.1969}}; {\citealt[]{Bolinger.1972}}; {\citealt[]{Os.1989}}; {\citealt[]{Kennedy.2005}}; {\citealt[]{Breindl.2007}}). Um das Zutreffen dieser Annahme zum ersten Mal mit empirischen Methoden zu testen, wird in den experimentellen Untersuchungen bspw. eine kontinuierliche Skala als {Bewertungsinstrument herangezogen}.


 

%% Wandel + Vielzahl
%\begin{singlespace*}
%\begin{quotation} "Das Bedürfnis nach Expressivität veranlasst die Sprecher, neue, stärkere, materialreichere Formen zu bilden und so die Entstehung immer neuer grammatischer Formen zu provozieren. Selbst wenn eine Sprache schon ein bestimmtes standardisiertes Ausdrucksmittel für eine grammatische Kategorie besitzt, können aufgrund des Bedürfnisses nach Expressivität neue, materialreiche Umschreibungen gefunden werden, die dann, wenn sie sich ebenfalls als Standardausdrücke durchsetzen, natürlich nicht mehr expressiv sind." \citep[106]{Diewald.1997}. \end{quotation}
%\end{singlespace*}

%Suscinskij (1985: 95): Vielzahl von Realisierungsmitteln + Auflistung: 1. steigernde Präfixe, 2. steigernde Basismorpheme, 3. steigernde Lexeme, 4. steigernde Wortgruppen, 5. Steigerungs-(Intensivierungs-)Sätze

%Berz (1953: 18): Wie mannigfach und groß das Bedürfnis nach expressiver Entspannung ist, das zeigt das lange Register an Auswahlmöglichkeiten, um ein und denselben Begriff zu verstärken. Da gibt es außer den schwachen Ausdrucksweisen mit \textit{ganz}, \textit{durchaus}, \textit{sehr}, den stärkeren mit \textit{furchtbar}, \textit{ungeheuer}, \textit{heidenmäßig} usw. und außer den ganz starken fluchenden \textit{verdammt}, \textit{verflucht} usw., allein innerhalb der zusammengesetzten Formen einen erstaunlich großen Reichtum an Abtönungen, die nicht Bedeutungsunterschiede zum Ausdruck bringen wollen, sondern feine, rational nur schwer zu erläuternde Abstufungen der Gefühlsbetontheit.

%Klara 2009: 35: Im Unterschied dazu weist ein Steigerungskompositum eine reichere kategoriale Füllung auf: zu der adjektivischen Basis können Nomina (\textit{honigsüß}, Adjektive \textit{hochmodern}, Verben \textit{kotzschlecht}, Präpositionen \textit{überglücklich} und Interjektionen \textit{pitschnass} als Steigerungsglieder hinzutreten.

%Klara 2009: 39: Abbildung 4: Intensivierungsmöglichkeiten im Adjektivbereich

%Setayesh & Vaez-Dalili (S. 50): interesting subject to study based on two charachteristics: fristly, because of their versatility and color, and secondly, because of their capacity for rapid change and recycling of forms. ... Intensification ist employed in both written and spoken language as a resource to convey a message in a more expressive way and to strengthen the speaker's position as well as their attitude toward what they are saying.

% Pittner in Klein, Duteil, Wagner 1991: 228: Der Wortbildungstyp STEIGERUNGSBILDUNG ist sehr produktiv. Zum einen können mit den meisten Steigerungsgliedern jederzeit neue Wörter gebildet werden, zum anderen kann auch nahezu jedes deutsche Wort als Steigerungsglied verwendet werden, d.h. die Klasse der Steigerungsglieder ist absolut offen. Trotzdem sind manche Steigerungsglieder offenbar nicht mehr produktiv, man denke nur an \textit{wunder}-, das zwar ältere Bildungen wie \textit{wunderschön} und \textit{wunderhübsch} zuläßt, sich aber nicht mit neueren Adjektiven vergleichbarer Semantik verbinden läßt \textit{*wundergeil, *wunderscharf, *wunderheiß}.






%% Inhaltliches
%Suscinskij (1985: 95): Vielzahl von Realisierungsmitteln + Auflistung: 1. steigernde Präfixe, 2. steigernde Basismorpheme, 3. steigernde Lexeme, 4. steigernde Wortgruppen, 5. Steigerungs-(Intensivierungs-)Sätze

%Suscinskij (1985: 96): Mittel, die den hohen Grad eines statischen bzw. dynamischen Merkmals anzugeben vermögen.

%Suscinskij (1985: 97): Potenziatoren neben dem Ausdruck eines hohen Grades auch pragmatische Funktionen ausüben, die vor allem in der Akzentuierung der Schreibereinstellung, in der Bekräftigung der vollen Übereinstimmung des Inhalts der Aussage mit der Wirklichkeit sowie in der bestimmten zielgerichteten Einwirkung auf den Textrezipienten.

%Suscinskij (1985: 97): Daraus resultiert, daß die Steigerung sowohl auf der Inhalts- als auch auf der Intentionsebene liegt. Die pragmatische Leistung der Potenziatoren ist es also, dazu beizutragen, daß der Hörer/Leser mit dem Sprecher/Schreiber gänzlich übereinstimmt, sein Urteil über einen bestimmten Sachverhalt akzeptiert und sich dementsprechend verhält, denkt bzw. handelt. Gerade in der Angabe der quantitativen Graduierung einer Eigenschaft oder der Intensität eines Prozesses, in der Bekräftigung des positiven Wahrheitswertes des Äußerungsinhalts oder im Streben des Sprechenden danach, den Empfänger in gewünschter Richtung zu beeinflussen, sind nach unserer Meinung die Ursachen für die Verwendung von Steigerungsmitteln in der Rede zu suchen.

%Suscinskij (1982: 329): Erhöhung der logisch-modalen Wirkung einer Aussage + Intensivierung der Aussageleistung + bei der Steigerung dominiert der logisch-semantische Aspekt und bei der Hervorhebung der kommunikativ-pragmatische

%Suscinskij (1981: 142): Die individuelle Beurteilung des Sachverhalts durch den Sprecher kann sich im Ausdruck von Emotionen [Befriedigung, Freude, Verachtung, Wut usw.], Wertung [positiv, negativ], Modalität [Sicherheit, Unsicherheit, Vermutung], Hervorhebung oder Abschwächung einer Aussage äußern.
%Suscinskij (1981: 142): Die logisch-emotionale Intensivierung einer Aussage -- die Hervorhebung -- wird durch bestimmte Lexik, bestimmte syntaktische und intonatorische Strukturschemata, durch eine besondere Satzgliedfolge, durch Wiederholung von Satzgliedern und durch Anschaulichkeit der Rede erreicht.

%Pustka (2014: 120): Expressivität im engeren Sinne [...] meint, der Etymologie entsprechend (vgl. lat. \textit{ex-pressus} 'ausgedrückt'), den Ausdruck von Emotionen. [...] Bei Expressivität im weiteren Sinne [...] steht dagegen die Darstellungsfunktion im Vordergrund

% Ruf 1996: 45: Die Verstärkung, Steigerung oder Intensivierung umfaßt also in einer weiten Auslegung des Begriffs alle Sprachmittel, die eine Intensivierung der Aussagekraft, eine Verstärkung des emotionalen Gehaltes oder der kommunikativen Absicht bewirken. Sie verleiht dem Ausdruck, Gewicht, Kraft und Bestätigung und bewirkt eine Nachdrücklichkeit der Rede.

% Ruf 1996: 46: Unter der Verstärkung im engeren Sinne kann man den Ausdruck des hohen Grades oder der hohen Intensität einer bezeichneten Eigenschaft oder eines Vorganges verstehen, sowie die Bezeichnung der besonderen Größe oder Ausdehnung des Genannten.

% Ruf 1996: 51: Intensivierungen drücken die Beziehung des Sprechers zum mitgeteilten Inhalt aus und sind somit eine subjektive Stellungnahme des Sprechers zum Gesagten

% Hentschel / Weydt in Weydt (1989:11): Die Bezeichnung Intensivpartikel lehnt sich an den englischen Begriff intensifier von Quirk/Greenbaum (1984: 244ff.) an. Intensivpartikeln haben die Funktion, die in einem anderen Wort enthaltene Eigenschaft zu modifizieren (zu verstärken oder abzuschwächen). Intensivpartikeln sind z. B. \textit{sehr, ziemlich, ganz} oder umgangsspprachlich \textit{irre, voll, echt}. Intensivpartikeln können Adjektive und zuweilen auch Verben als Bezugswörter haben: \textit{Ich bin ziemlich/sehr müde, Sie hat mich ziemlich/sehr geärgert}. Sie müssen mit ihrem Bezugswort zusammen erfragt werden und können nur mit diesem zusammen die Antwort auf eine Frage bilden. Intensivpartikeln lassen sich semantisch in verstärkende (\textit{sehr, höchst, irre}, neu: voll) und abschwächende (\textit{ziemlich, etwas, einigermaßen}) einteilen. Mit Ausnahme von \textit{sehr} und \textit{besonders} können bei Intensivpartikeln normalerweise keine Negationen stehen (vgl. \textit{Nicht sehr/besonders nett} gegenüber \textit{*nicht ziemlich traurig, *nicht höchst traurig}.

% Wolski in Weydt (1989: 350): Die propositionale Bedeutung gilt als der bewertbare Teil der Satzbedeutung. Die nicht-propositionale Bedeutung als der bewertende Teil der Satzbedeutung. Durch letzteren werden jemandes Einstellungen zu dem, was mit der propositionalen Bedeutung identifiziert wird, spezifiziert. Der bewertende Teil kann epistemischer, intentionaler und emotionaler Art sein.

%erstes Element, das semantisch häufig gänzlich ausgebleicht ist, verstärkt das zweite (vlg. Pittner in Klein, Duteil, Wagner 1991: 225)




%% Arten
% 2 Möglichkeiten: a) deskr. Gehalt + Emotionalität = lexikalisierte Formen wie Gaul, Köter etc.; b) Emotionalität + deskriptiver Gehalt = expressive Konstruktion, bei der sich die Emotionalität aus der Konstruktion heraus ergibt (sau langweilig)

% Allgemeine Fragestellungen: Werden Adjektive, die als Intensivierer fungieren, noch als emotional wahrgenommen? Oder nur als graduierend? Wo auf der Skala verortet? Wie sehr verstärkt der Intensivierer die Emotionalität? Wie werden Intensivierer hinsichtlich Intensität und Emotionalität vewertet (Intensität skalierbar, Emotionalität skalierbar, aber beides gleich?)?



%% Frequenz/Zeit
%Richter & van Hout (2020: 343): The most evident hypothesis is that an established intensifier has a high use value, but a low expressive value. Non-established intensifiers will have low use values and high expressive values.

%Richter & van Hout (2020: 354): modern intensifiers may cause surprisal and attract attention: they may unfold a high intensifying effect while still carrying a great deal of the original meaning



%% Abgrenzungen
%lexikalische Graduierung versus morphologische oder syntaktische (vgl. Metzler Lexikon Sprache 2000)

%Kirschbaum 2002: 201: Morphologische versus syntaktische Intensivierung --> wir untersuchen syntaktische, d. h. getrennt: sau heiß und nicht sauheiß




%% Sonstiges:
%Terminus expressiv stammt aus der Musik

% was im Rahmen der Arbeit nicht gemacht wird: der morphologische Status von Intensivierern diskutieren (s. Klara; Pittner, Robert usw)

% Hauschild (1899: 1): mehr der Sprache des täglichen Lebens angehörend







\section{Zielsetzung und Forschungshypothesen}\label{abs:ziel}
Der Fokus dieses Buchs liegt auf der experimentellen Untersuchung expressiver Intensitätspartikeln wie \textit{mega} oder \textit{super} und der Einschätzung der durch sie ausgedrückten Intensität. Dies dient der Überprüfung zentraler Annahmen, die in der theoretischen Literatur zwar weithin akzeptiert werden, deren Gültigkeit bisher allerdings nicht abgesichert wurde. Eine solche empirische Vorgehensweise stellt daher bis dato ein Desiderat dar. In dieser Studie wird auf eine umfassende Kombination aus verschiedenen linguistischen Standardmethoden wie Korpusanalyse und experimentelle Untersuchungen zurückgegriffen, mithilfe derer die hergeleiteten Forschungshypothesen entweder bestätigt oder entkräftet werden.

In der in \chapref{Korpus} präsentierten Korpusstudie wird eine Häufigkeitsanalyse ausgewählter expressiver Intensivierer (n = 16) des Gegenwartsdeutschen\footnote{Auch wenn die Sprachperiodisierung des Deutschen mitnichten einheitlich definiert ist, ist das Gegenwartsdeutsche in etwa in die Zeit ab Mitte des 20. Jahrhunderts zu verorten \citep[vgl.][3]{Elspass.2008}.} vorgenommen, mit dem Gedanken die synchrone und diachrone Entwicklung ebendieser in einem auf die Ziele der Studie ausgerichteten Textkorpus sichtbar zu machen. Im Mittelpunkt stehen die folgenden Hypothesen:

\begin{itemize}
\item[\textbf{\textsc{H1}}] Eine hochfrequente Verwendung von Intensivierern führt unweigerlich zum Verlust ihrer Expressivität.

\item[\textbf{\textsc{H2}}] Der Expressivitätsverlust geht mit einem Bedarf an neuen Ausdrücken einher, mit denen sich Sprecher auch weiterhin als expressiv, innovativ und originell darstellen können.
\end{itemize}

\noindent In diesem Sinne liegt das Ziel der Korpusstudie einerseits in der Beantwortung der Frage, ob sich ein frequenter Gebrauch de facto nachteilig auf die Expressivität von Intensivierern auswirkt, wie neben \citet[]{Hauschild.1899} und {\citet[]{Keller.1995}} auch \citet[]{Robertson.1954} bemerkt:\footnote{Da eine frequenzbasierte Korpusanalyse nur das zeitliche Auftreten von Ausdrücken in einer Textsammlung erfasst, können die gewonnenen Daten zwangsläufig keine direkten Aussagen über den Grad der Expressivität der Intensivierer machen. Dennoch bieten sie erste Informationen in die Zusammenhänge von Frequenz und Persistenz und können als Indikator für einen möglichen Verlust an Expressivität dienen. Die longitudinale Entwicklung wird daher als erster Hinweis auf den Expressivitätsgrad jedes Intensivierers genutzt.}

% \begin{singlespace*}
\begin{quotation} When the strong word is used on light occasion its strength begins to be dissipated, and when the fitting moment for it actually arrives it will no longer serve; familiarity has bred contempt in the hearer, and one must begin again to find a new `strong word.' \hfill \citep[251]{Robertson.1954} \\ \end{quotation} 
% \end{singlespace*} 

\noindent Dieser Aspekt wird auch in einigen neueren Arbeiten behandelt (bspw. \citealt[]{Bolinger.1972}; \citealt[]{Ito.2003}; \citealt[]{Tagliamonte.2008}; {\citealt[]{Stratton.2020d,Stratton.2022a}}). So nimmt \citet[18]{Bolinger.1972} bspw. an, dass die Wahl eines Intensivierers vom Wunsch des Sprechers motiviert ist, originell zu sein, und in einer ständigen sprachlichen Erneuerung Ausdruck findet: \glqq{}Degree words afford a picture of fevered invention and competition that would be hard to come by elsewhere, for in their nature they are unsettled.\grqq{} Andererseits wird auch die Aussage, dass die pragmatische Abnutzung in der Entwicklung neuer Ausdrücke resultiert, durch die Untersuchung hinterfragt. Diese geht in der Literatur vornehmlich auf {\citet[]{Becher.1907}} und {\citet[]{Baumgarten.1908}} zurück. Was die {Hypothese H1} neben dem Faktor Frequenz ebenso zu beinhalten scheint, ist der Faktor Zeit. So ist zu bedenken, dass die postulierte Abnutzung nicht aus einem Punkt heraus geschieht, sondern sich über einen gewissen Zeitraum erstreckt, wie bereits {\citet[]{Tobler.1868}} und {\citet[]{Gabelentz.2016}} erkennen. Mit anderen Worten: Je länger eine Intensitätspartikel im Gebrauch ist, desto weniger überraschend und aktuell ist sie, weshalb sie gleichsam mit einem niedrigeren Grad an Expressivität assoziiert wird: \glqq{}The most evident hypothesis is that an established intensifier has a high use value, but a low expressive value. Non-established intensifiers will have low use values and high expressive values\grqq{} \citep[346f.]{Richter.2020}. Die intensivierende Wirkung lässt also im Laufe der Zeit und mit zunehmender Häufigkeit nach: \glqq{}Overuse, diffused use, long-time use, will lead to a diminishment in the intensifier’s ability to boost and intensify\grqq{} ({\citealt[391]{Tagliamonte.2008}}). Deshalb gilt es, zusätzlich zur Verwendungshäufigkeit die Zeit bzw. das Erstvorkommen im Korpus hinzuzuziehen, weswegen die Hypothese H1 wie folgt spezifiziert wird:

\begin{itemize} 
\setlength{\itemindent}{1cm}
\item[\textbf{\textsc{H1-Freq}}] Eine hochfrequente Verwendung von Intensivierern führt zu einer pragmatischen Abnutzung und demzufolge sukzessive zum Verlust von Expressivität. 

\item[\textbf{\textsc{H1-Temp}}] Eine persistente Verwendung von Intensivierern führt zu einer pragmatischen Abnutzung und demzufolge sukzessive zum Verlust von Expressivität. 
\end{itemize}

Da es durch eine quantitative Korpusanalyse nicht möglich ist, stichfeste Aussagen über den tatsächlichen Expressivitätsstatus der Ausdrücke zu machen, lassen sich damit doch lediglich Entwicklungstendenzen ableiten, wird die Beantwortung dieser Fragen auf Grundlage der erzielten Ergebnisse anschließend mittels experimenteller Methoden angestrebt. Diese in {\chapref{Experiment0}} geschilderte Untersuchung erfolgt in Gestalt von Evaluationsaufgaben, bei denen die in der Korpusstudie berücksichtigten Intensivierer von einer Stichprobe an Versuchspersonen stellvertretend für die Grundgesamtheit bewertet wird. Die Bewertung wird dabei im Hinblick auf Skalarität und Sprechereinstellung vorgenommen, wie in folgenden Hypothesen {festgehalten ist}:

\begin{itemize}
\item[\textbf{\textsc{H3}}] Expressive Intensivierer werden mit einem höheren Grad an Skalarität
assoziiert als deskriptive Intensivierer. 

\item[\textbf{\textsc{H4}}] Expressive Intensivierer geben eine Einstellungsbekundung des Sprechers
zu dem im Satz beschriebenen Sachverhalt wieder und werden in der Folge mit einem höheren Grad an Sprechereinstellung assoziiert als deskriptive Intensivierer. 
\end{itemize}

\noindent In der theoretischen Literatur wird also davon ausgegangen, dass expressive Intensivierer zum einen ein höheres Maß an Merkmalsausprägung angeben und somit in Bezug auf Skalarität auf einer Skala eine höhere Position einnehmen als deskriptive Intensivierer. Diese Einschätzung findet sich bspw. in {\citet[]{Paradis.1997}} und \citet[]{Gutzmann.2019a}. Zum anderen spiegeln expressive Intensivierer auch die persönliche Einstellung des Sprechers wider, wodurch sie ebenfalls auf einer Skala der kommunizierten Sprechereinstellung höher verortet werden. Diese These geht u.\,a. auf \citet[]{Lang.1983} und {\citet[]{Ortner.2014}} zurück. Die in diesem Kapitel vorgestellte experimentelle Untersuchung verfolgt daher in erster Linie das Ziel, beide Thesen zum ersten Mal auf ihre Korrektheit hin zu testen. Daneben wird die Überprüfung einer weiteren in diesem Kontext relevanten Hypothese angestrebt, die sich aus einer Schlussfolgerung von \citet[]{Gutzmann.2019a} heraus ergibt, der annimmt, dass die ausgedrückte Sprechereinstellung evaluativ ist und sich in Form einer positiven oder negativen Bewertung des im Satz thematisierten Sachverhalts niederschlägt. Im Gegensatz dazu dienen deskriptive Intensivierer vermutlich nicht der Einstellungsspiegelung und werden bezüglich der kommunizierten Sprechereinstellung vorwiegend als neutral beurteilt. Die zu überprüfende Hypothese lautet deswegen wie folgt:

\begin{itemize} 
\item[\textbf{\textsc{H5}}] Die Einstellungsbekundung, die expressive Intensivierer zu dem im Satz beschriebenen Sachverhalt wiedergeben, ist evaluativer Natur und wird daher als positiv oder negativ wahrgenommen.
\end{itemize}

Im Fokus der in \chapref{Experimente} beschriebenen Untersuchung steht eine weiterführende Hypothese, die sich implizit aus der Annahme eines wahrnehmbaren Unterschiedes zwischen deskriptiven und expressiven Intensivierern ergibt. So lässt sich auf Basis der theoretischen Literatur (bspw. {\citealt[]{Biedermann.1969}} oder \citealt[]{Breindl.2007}) erwägen, dass sich nicht nur die beiden Intensiviererklassen hinsichtlich Skalarität und Sprechereinstellung voneinander unterscheiden, sondern sich auch klasseninterne Abstufungen zeigen, die semantisch bedingt sind. Demnach lassen sich die Mitglieder einer Klasse mutmaßlich in Abhängigkeit des durch sie ausgedrückten Grades an Skalarität und Sprechereinstellung auf einer kontinuierlichen Skala längs anordnen. Um die vermuteten Abstufungen empirisch zu überprüfen, wurden Evaluationsaufgaben erstellt, bei denen eine Personenstichprobe jeweils vier deskriptive und expressive Intensivierer in Bezug auf beide Merkmale beurteilte. Die Untersuchung basiert auf folgender Hypothese:

\begin{itemize}
\item[\textbf{\textsc{H6}}] Innerhalb der Intensiviererklassen lassen sich sowohl im Hinblick auf Skalarität als auch auf Sprechereinstellung graduelle Abstufungen zwischen den Ausdrücken beobachten, die semantisch bedingt sind.
\end{itemize}

\noindent Dabei ist im Sinne von \citet[]{Gutzmann.2019a} anzunehmen, dass die Expressivität eines Intensivierers zu einer Anhebung des durch ihn ausgedrückten Skalaritätsgrades führt, der in der Konsequenz als stärker wahrgenommen wird. In dem Fall würden die beiden Merkmale korrelieren. Ebendiesen potenziellen Zusammenhang zu untersuchen, ist eines der Ziele dieses Experiments.

Zusammengefasst wird sich im Rahmen der Untersuchung mit Korpusstudie und Experimenten einer systematischen Kombination aus verschiedenen Methoden bedient, mithilfe derer erstmals belastbare Aussagen über die Gültigkeit der zur Diskussion stehenden Hypothesen gemacht werden können. 







%Da die Hypothesen inhaltlich aufeinander aufbauen, werden sie in diesem Abschnitt nicht in der Reihenfolge vorgestellt, in der sie in der Arbeit Beachtung finden. Stattdessen werden sie hier im Sinne der Nachvollziehbarkeit angeordnet. Begonnen wird mit den Forschungshypothesen, deren Überprüfung in der in \chapref{Experimente} präsentierten Untersuchung angestrebt wird.

%Im Fokus der in \chapref{Experimente} dargelegten Experimente stehen Evaluationsaufgaben, bei denen eine Personenstichprobe jeweils vier deskriptive und expressive Intensivierer des Gegenwartsdeutschen\footnote{Auch wenn die Sprachperiodisierung des Deutschen mitnichten einheitlich definiert ist, ist das Gegenwartsdeutsche in etwa in die Zeit ab Mitte des 20. Jahrhunderts zu verorten \citep[vgl.][3]{Elspass.2008}.} stellvertretend für die Grundgesamtheit bewertet. Dies geschieht vor dem Hintergrund der oben beschriebenen Annahme, dass es einen grundlegenden Unterschied zwischen der Wahrnehmung von deskriptiven und expressiven Intensivierern gibt, der sich auf eine Skala abbilden lässt. Dementsprechend fußt die Untersuchung auf den folgenden beiden Hypothesen:
%
%\begin{itemize}
%\item[\textbf{\textsc{h4}}] Expressive Intensivierer werden mit einem höheren Grad an Skalarität
%assoziiert als deskriptive Intensivierer. 
%
%\item[\textbf{\textsc{h5}}] Expressive Intensivierer geben eine Einstellungsbekundung des Sprechers
%zu dem im Satz beschriebenen Sachverhalt wieder und werden in der Folge mit einem höheren Grad an Sprechereinstellung assoziiert als deskriptive Intensivierer. 
%\end{itemize}
%
%Es wird also davon ausgegangen, dass expressive Intensivierer zum einen einen höheren Grad einer Merkmalsausprägung angeben und somit in Bezug auf Skalarität auf einer Skala eine höhere Position einnehmen als deskriptive Intensivierer. Diese Behauptung findet sich bspw. in \citet[]{Paradis.1997} und \citet[]{Gutzmann.2019a}. Zum anderen spiegeln expressive Intensivierer auch die persönliche Einstellung des Sprechers wider, wodurch sie auf einer Skala der kommunizierten Sprechereinstellung gleichfalls höher verortet werden als deskriptive Intensivierer. Diese These geht in der theoretischen Literatur u.\,a. auf \citet[]{Lang.1983} und \citet[]{Ortner.2014} zurück. Die in diesem Kapitel vorgestellte experimentelle Untersuchung verfolgt daher in erster Linie das Ziel, beide Thesen zum ersten Mal auf ihre Gültigkeit hin zu überprüfen.

%Daneben wird in der in \chapref{Korpus} präsentierten Korpusstudie eine Häufigkeitsanalyse ausgewählter expressiver Intensitätspartikeln (n = 16) vorgenommen, mit dem Gedanken, die synchrone und diachrone Entwicklung ebendieser in einem auf die Ziele der Arbeit ausgerichteten Textkorpus sichtbar zu machen. Im Mittelpunkt stehen dabei die folgenden Hypothesen:
%
%\begin{itemize}
%\item[\textbf{\textsc{h1}}] Eine hochfrequente Verwendung von Intensivierern führt unweigerlich zum Verlust ihrer Expressivität.
%
%\item[\textbf{\textsc{h2}}] Der Expressivitätsverlust geht mit einem Bedarf an neuen Ausdrücken einher, mit denen sich Sprecher auch weiterhin als expressiv, innovativ und originell darstellen können.
%\end{itemize}
%
%In diesem Sinne liegt das Ziel der Korpusstudie einerseits in der Beantwortung der Frage, ob sich ein frequenter Gebrauch \textit{de facto} nachteilig auf die Expressivität von Intensivierern auswirkt, wie u.\,a. \citet[]{Hauschild.1899} und {\citet[]{Keller.1995}} behaupten. Andererseits wird auch die Aussage, dass die pragmatische Abnutzung mit der Entwicklung neuer Ausdrücke einhergeht, durch die Untersuchung hinterfragt. Diese geht in der Literatur primär auf {\citet[]{Becher.1907}} und {\citet[]{Baumgarten.1908}} zurück. Was H1 neben dem Faktor Frequenz ebenso zu beinhalten scheint, ist der Faktor Zeit. So ist davon auszugehen, dass die postulierte Abnutzung nicht aus einem Punkt heraus geschieht, sondern sich über einen gewissen Zeitraum erstreckt, wie bereits \citet[]{Tobler.1868} erkennt. Mit anderen Worten: Je länger eine Intensitätspartikel im Gebrauch ist, desto weniger überraschend und aktuell ist sie, wodurch sie gleichsam mit einem niedrigeren Grad an Expressivität assoziiert wird: "The most evident hypothesis is that an established intensifier has a high use value, but a low expressive value. Non-established intensifiers will have low use values and high expressive values" \citep[346f.]{Richter.2020}. In der Folge gilt es zusätzlich zur Verwendungshäufigkeit die Zeit bzw. das Erstvorkommen im Korpus hinzuzuziehen, weswegen H1 wie folgt {spezifiziert wird}:
%
%\begin{itemize} 
%\setlength{\itemindent}{1cm}
%\item[\textbf{\textsc{h1-freq}}] Eine hochfrequente Verwendung von Intensivierern führt zu einer pragmatischen Abnutzung und demzufolge sukzessive zum Verlust von Expressivität. 
%
%\item[\textbf{\textsc{h1-temp}}] Eine persistente Verwendung von Intensivierern führt zu einer pragmatischen Abnutzung und demzufolge sukzessive zum Verlust von Expressivität. 
%\end{itemize}
%
%Da es mithilfe einer quantitativen Korpusanalyse nicht möglich ist, stichfeste Aussagen über den tatsächlichen Expressivitätsstatus der Ausdrücke zu machen, lassen sich damit doch lediglich Entwicklungstendenzen ableiten, wird die Beantwortung dieser Fragen auf Grundlage der erzielten Ergebnisse anschließend mittels experimenteller Methoden angestrebt. Diese in {\chapref{Experiment0}} geschilderte Untersuchung erfolgt in Gestalt von Evaluationsaufgaben, bei denen die in der Korpusstudie berücksichtigten Intensivierer von einer Stichprobe an Versuchspersonen im Hinblick auf Skalarität (vgl. H4) und Sprechereinstellung (vgl. H5) bewertet wird. Ergänzend dazu wird untersucht, ob expressive Intensivierer neben Skalarität auch die evaluative Haltung des Sprechers zum Geäußerten ausdrücken. Wie im Hinblick auf H5 skizziert, vertreten bspw. {\citet[315]{Lang.1983}} und \citet[57]{Ortner.2014} die Auffassung dass, expressive Intensivierer eine Einstellungsbekundung wiedergeben und infolgedessen mit einem höheren Grad an Sprechereinstellung assoziiert werden als deskriptive Intensivierer. Vor diesem Hintergrund ist im Sinne von \citet[681]{McCready.2008} und \citet[135]{Gutzmann.2019a} anzunehmen, dass die ausgedrückte Sprechereinstellung evaluativer Natur ist und sich in Form einer positiven oder negativen Bewertung des beschriebenen Sachverhalts niederschlägt. Deskriptive Intensivierer geben dagegen mutmaßlich keine Einstellungsbekundung wieder und werden demnach in Bezug auf die kommunizierte Sprechereinstellung vorwiegend als 'neutral' beurteilt. Die zu testende Hypothese lautet daher {wie folgt}:
%
%\begin{itemize} 
%\item[\textbf{\textsc{H3}}] Die Einstellungsbekundung, die expressive Intensivierer zu dem im Satz beschriebenen Sachverhalt wiedergeben, ist evaluativer Natur und wird daher als positiv oder negativ wahrgenommen.
%\end{itemize}
%
%Zusammengefasst wird sich im Rahmen der Arbeit mit Korpusstudie und Experimenten einer systematischen Kombination aus verschiedenen Methoden bedient, mithilfe derer erstmals belastbare Aussagen über die Gültigkeit der zur Diskussion stehenden Hypothesen gemacht werden können. 


\section{Kapitelgliederung}\label{abs:gliederung}
Nachdem dieses Kapitel den Problembereich, die Motivation und die Forschungshypothesen definiert hat, werden in \chapref{Theorie} der theoretische Rahmen gesetzt und thematisch relevante Konzepte vorgestellt. Im ersten Teil erfolgt ein Forschungsüberblick, in dem die Untersuchung in den Kontext früherer Forschungsbeiträge eingegliedert und ihre wissenschaftliche Relevanz begründet wird. Daraufhin werden elementare Begriffe geklärt, die im Buch verwendet werden. Dieses Unterkapitel zerfällt in zwei Teile: Zum einen wird der Unterschied zwischen deskriptivem und expressivem sprachlichen Gehalt herausgearbeitet, zum anderen erfolgt eine Analyse der beiden Klassen von Intensitätspartikeln, die den zentralen {Untersuchungsgegenstand bilden}.

An den theoretischen Hintergrund schließt sich in \chapref{Korpus} der erste empirische Teil an, dessen Kern eine frequenzbasierte Auswertung von Intensivierern ist. Die darin präsentierte Korpusstudie verfolgt das Ziel, ausgewählte Intensivierer des Gegenwartsdeutschen synchron und diachron über einen Zeitraum von etwa 70 Jahren quantitativ einzuordnen, um stichhaltige Hinweise auf deren Gebrauch in einem schriftsprachlichen Kontext zu erlangen. Das Kapitel setzt sich ganz klassisch aus theoretischem Rahmen, methodischem Vorgehen, Ergebnispräsentation, Diskussion und Fazit zusammen. Die Ergebnisse der Korpusanalyse werden im weiteren Verlauf als Grundstein für die darauffolgenden experimentellen Untersuchungen genutzt, die den zweiten Empirieteil des Buchs darstellen.

Bevor die Experimente im Detail beschrieben werden, sei an dieser Stelle darauf hingewiesen, dass ihre Anordnung im Buch aus Gründen der Nachvollziehbarkeit einer logischen Darstellung folgt, während die tatsächliche zeitliche Abfolge davon abweicht: In der Praxis wurde zunächst das in {\chapref{Experimente}} präsentierte Experiment 2 durchgeführt, an das sich das in {\chapref{Experiment0}} dargelegte Experiment 1 anschloss. Die Darstellung der Experimente im Buch erfolgt jedoch in antichronologischer Reihenfolge. Motiviert ist die zeitliche Umkehr zum einen durch den Untersuchungsgegenstand, der in der Korpusstudie und Experiment 1 identisch ist. Um weitere Informationen zu den nach der Korpusanalyse noch offenen Fragestellungen zu gewinnen, wurden die korpuslinguistisch überprüften Intensivierer ergänzend experimentell mittels Skalenrating untersucht. Zum anderen bauen die Korpusstudie und Experiment 1 auch inhaltlich in Bezug auf die zentralen Forschungshypothesen aufeinander auf, was ebenfalls für eine Abkehr von der tatsächlichen zeitlichen Reihenfolge spricht.

\chapref{Experiment0} stellt die erste experimentelle Untersuchung vor, bei der die in der Korpusstudie berücksichtigten expressiven Intensivierer in Form einer Rating-Studie relativ zueinander in Beziehung gesetzt werden. Dieses Brückenexperiment dient der weiteren Überprüfung der der Korpusanalyse zugrunde gelegten Hypothesen. Auch hier wird erst der theoretische Rahmen abgesteckt, an den sich die zu testenden Forschungshypothesen, der Untersuchungsaufbau und Untersuchungsgegenstand sowie Versuchspersonen und Auswertungskriterien anschließen. Darauf folgt die Ergebnispräsentation, die sich wiederum aus einer deskriptiv- und inferenzstatistischen Datenauswertung zusammensetzt. Diese Ergebnisse werden hinterher resümiert und interpretiert, woraufhin die fokussierten Hypothesen im Fazit einer kritischen Betrachtung unterzogen werden.

\chapref{Experimente} widmet sich der zweiten experimentellen Untersuchung, in der eine Stichprobe von jeweils vier deskriptiven und expressiven Intensivierern mithilfe eines Skalenratings hinsichtlich Skalarität und Sprechereinstellung miteinander verglichen wird. Wie das vorangegangene Kapitel besteht auch dieses zunächst aus theoretischem Rahmen und Forschungshypothese nebst -aufbau. Der Untersuchungsteil ist in zwei separate Segmente gegliedert: Während sich das erste Segment der Untersuchung von Skalarität widmet, ist im zweiten Segment die Sprechereinstellung vordergründig. Vor diesem Hintergrund werden jeweils Untersuchungsgegenstand und -durchführung, Versuchspersonen und Auswertungskriterien vorgestellt. Wiederholt setzt sich die Ergebnispräsentation aus einer deskriptivstatistischen und inferenzstatistischen Datenanalyse zusammen, an die sich eine erste Diskussion der Ergebnisse anschließt. Daraufhin wird dem möglichen statistischen Zusammenhang zwischen Skalarität und Sprechereinstellung nachgegangen. Abgeschlossen wird das Kapitel durch eine Gesamtdiskussion der untersuchten Fragestellungen sowie ein Fazit, in dem die Forschungshypothese kumuliert für beide Experimente auf ihr Zutreffen hin inspiziert wird.

Das sich daran anschließende \chapref{Abschlussdiskussion} fasst die Ergebnisse des empirischen Teils unter Bezugnahme der zu überprüfenden Hypothesen in einer abschließenden Diskussion final zusammen.

 Schließlich werden in \chapref{Gesamtfazit} alle erzielten Erkenntnisse resümiert und ein Ausblick auf offene Forschungsfragen gegeben, die im Rahmen dieser Studie nicht oder bloß unzureichend beantwortet wurden und Gegenstand nachfolgender Untersuchungen auf diesem Gebiet sein können.





